%!TeX program = xelatex
\documentclass{SYSUReport}
\usepackage{amsmath}
\usepackage{tikz}
\usepackage{tikz-qtree}
\usepackage{adjustbox} % 加入调整表格大小的宏包
\usepackage{booktabs} % 用于创建更漂亮的表格
\usepackage{multirow} % 用于合并单元格
\usepackage{graphicx}
\usepackage{longtable}
\usepackage{amsmath}

\usepackage{listings}
\usepackage{xcolor}

\lstdefinestyle{mystyle}{
    backgroundcolor=\color{gray!10},      % 背景色
    basicstyle=\ttfamily\small,           % 基本字体
    keywordstyle=\color{blue!70!black},   % 关键字颜色
    commentstyle=\color{green!40!black},  % 注释颜色
    stringstyle=\color{purple!60!black},  % 字符串颜色
    numbers=left,                         % 行号
    numberstyle=\tiny\color{gray},
    frame=single,                         % 边框
    breaklines=true,                      % 自动换行
    captionpos=b                          % caption 在底部
}

\lstset{style=mystyle}


% 根据个人情况修改
\headl{学号 姓名}
\headr{数据安全和隐私保护}
\lessonTitle{数据安全和隐私保护}
\reportTitle{数据安全和隐私保护第x次实验}
\stuname{姓名}
\stuid{学号}
\inst{网络空间安全学院}
\major{网络空间安全}
% 修改为具体日期
\date{2026 年 9 月 10 日}

\begin{document}

\cover
\thispagestyle{empty} % 首页不显示页码
\clearpage


\tableofcontents
\clearpage


\pagenumbering{arabic}
\setcounter{page}{1}

\section{实验要求}

把每次实验要求粘贴过来就行，不粘贴删掉这个section也行。

\section{实验原理}

根据具体实验简单介绍实验原理。

\section{实验设置}

简单介绍下实验使用的软硬件版本和复现方法/算法是什么。

\section{实验过程}

结合实验操作和实验结果描述实验过程，并附带必要的实验分析。

\section{实验总结}

写自己的实验总结、实验心得

%以下是一些贴图、引用和贴代码（示例为matlab）的示例，latex具体语法可以自行上网/gpt搜索
% 首先我们先对S-UNIWARD算法隐写的结果进行可视化，下图\ref{fig:bbp01}-图\ref{fig:bbp02}分别对应嵌入率payload=0.1、0.2、0.3、0.4、0.5bbp的结果。图中的cover为原图像，embedding changes对应嵌入信息，Stego Image对应隐写后图像。可以看到，随着嵌入率bbp的提高，嵌入的信息越多，对原图像的修改量就越大（但用肉眼仍然难以分辨）。

% \begin{figure}[h!t!b!p]
%     \centering
%     \includegraphics[width=0.8\linewidth]{bbp1.png}
%     \caption{payload=0.1bbp时隐写结果}
%     \label{fig:bbp01}
% \end{figure}

% \begin{figure}[h!t!b!p]
%     \centering
%     \includegraphics[width=0.8\linewidth]{bbp2.png}
%     \caption{payload=0.2bbp时隐写结果}
%     \label{fig:bbp02}
% \end{figure}

% \begin{lstlisting}[language=Matlab, caption={S-UNIWARD 隐写嵌入代码}]
% clc; clear all;

% coverPath = 'Bossbase_1.01/';
% payloads = [0.1, 0.2, 0.3, 0.4, 0.5];
% sPaths = {'stego_01','stego_02','stego_03','stego_04','stego_05'};
% \end{lstlisting}

% 如需参考文献可以取消以下两行代码注释并在reference.bib中添加bib引用，然后在正文中用\ref引用
% \bibliographystyle{ieeetr}   % 或 plain / unsrt / IEEEtran 等
% \bibliography{reference}    

\end{document}